\section{实践：Lab 2 — 面部检测系统}

课程配套了一个动手实验（Lab 2），使用 VAE 对人脸数据集进行训练，自动发现隐变量（如肤色、姿态、光照等），并利用这些隐变量检测数据集中的偏差和不足代表的特征。

\begin{figure}[H]
\centering
\includegraphics[width=0.85\textwidth]{slides-images/slide-89.jpg}
\caption{Lab 2：基于 VAE 的面部检测系统\protect\footnotemark}
\end{figure}
\footnotetext{Slides 第89页。}

\begin{knowledgebox}{Lab 2 的教学意义}
这个实验将理论与实践相结合：
\begin{itemize}
    \item 构建 VAE 模型处理人脸图像
    \item 通过隐变量扰动发现数据中的潜在特征
    \item 检测训练数据中的偏差（如某些肤色或姿态是否被低估）
    \item 利用生成模型创建更公平、更平衡的数据集
\end{itemize}
实验代码可在 \href{https://github.com/MITDeepLearning/introtodeeplearning}{github.com/MITDeepLearning/introtodeeplearning} 获取。
\end{knowledgebox}

\subsection{本章小结}

配套实验帮助学生在动手实践中深化对 VAE 的理解，特别是如何利用生成模型发现和缓解数据集中的偏差问题——这是 AI 公平性研究的重要课题。


